\chapter{Conclusion}
\label{chap:Conclusion}

À l'issue de cette première phase, nous résumons le travail effectué, en tirons les principaux enseignements et présentons les perspectives de la phase~II.

%--------------------------------------------------------------------------------

\section{Résumé du travail effectué}
\label{sec:Resume}

Nous avons étudié l'apport d'un transformeur de graphes 3D pré-entraîné, Uni-Mol, à la prédiction multi-label des 138 descripteurs olfactifs du jeu GS-LF.
L'étude bibliographique a situé ce choix par rapport aux GNN de référence, POM et OpenPOM, limités par la myopie de la propagation de messages et l'absence de géométrie, et nous a fait découvrir un travail récent appliquant Uni-Mol2 à la même tâche, dont nous avons repris le protocole d'évaluation.

Nous avons développé une chaîne d'apprentissage complète~: génération de conformations 3D, encodage au format Uni-Mol, augmentation géométrique, décodeur MLP, perte asymétrique et dégel partiel de l'encodeur.
Dix-sept expériences ont porté le F1 macro de validation de 0,205 à 0,354 et l'AUROC macro à 0,872.
Ces valeurs sont du même ordre que celles des modèles de référence (F1 de 0,375 pour OpenPOM, 0,399 pour Uni-Mol2 ajusté), mais ne leur sont pas encore comparables~: elles sont mesurées sur la validation, sans test tenu à l'écart.
Le protocole de référence, implémenté en fin de phase, permettra cette comparaison.

%--------------------------------------------------------------------------------

\section{Enseignements}
\label{sec:Enseignements}

\paragraph{Le protocole d'évaluation fait partie du résultat.}
Un F1 n'a de sens qu'avec le découpage, le choix des seuils et la sélection du modèle qui l'accompagnent.
Nous avons longtemps optimisé une mesure optimiste~; fixer le protocole de la littérature dès le début aurait permis des comparaisons fiables tout au long du projet.
De même, F1 micro et F1 macro mesurent des choses différentes~: le premier, dominé par les descripteurs fréquents, est bien plus élevé que le second.

\paragraph{Un modèle pré-entraîné exige un encodage exact.}
Les écarts d'encodage que nous avons corrigés ne produisaient aucune erreur, seulement des performances dégradées.
La correction a apporté l'un des gains les plus importants du projet~: vérifier la conformité des entrées au pré-entraînement est une étape à part entière.

\paragraph{Le déséquilibre se traite à plusieurs niveaux.}
La pondération des positifs, puis la perte asymétrique, ont apporté des gains nets, et le choix des seuils de décision en apporte d'autres.
Ces leviers agissent à des niveaux différents (apprentissage, décision) et se combinent.

\paragraph{Peu de données, beaucoup de paramètres.}
Le dégel de l'encodeur améliore les performances mais provoque un surapprentissage rapide~: sur 4\,000 molécules, chaque million de paramètres libérés doit être justifié.

\paragraph{Sur le plan de l'organisation.}
Le suivi systématique des expériences (configuration, courbes et commentaire pour chacune) a rendu ce rapport possible et permet de revenir sur des conclusions antérieures.
À l'inverse, modifier plusieurs paramètres à la fois a empêché d'attribuer certains gains.
\acompleter{enseignements personnels de chaque auteur}

%--------------------------------------------------------------------------------

\section{Perspectives de recherche}
\label{sec:Perspectives}

\paragraph{Évaluer selon le protocole de référence.}
La première étape de la phase~II est d'exécuter le protocole sur le jeu complet, afin de situer réellement notre modèle par rapport à POM, OpenPOM et Uni-Mol2.
Il faudra ensuite maîtriser le surapprentissage lié au dégel (taux d'apprentissage plus faible pour l'encodeur, sélection du modèle sur l'AUROC, dégel progressif), revérifier les conclusions des expériences antérieures à la correction de l'encodage, et mener une recherche d'hyperparamètres.

\paragraph{Passer à Uni-Mol2.}
Le sujet désigne Uni-Mol2 comme modèle cible.
Sa représentation de paires, qui combine géométrie et topologie, et sa taille supérieure devraient améliorer les performances, comme le suggère \cite{Han-2026}, au prix d'un coût de calcul plus élevé.

\paragraph{Modéliser les dépendances entre descripteurs.}
Le décodeur traite les 138 descripteurs comme indépendants~: chaque ligne $\mathbf{w}_j$ de la dernière matrice $\mathbf{W} \in \mathbb{R}^{138 \times 256}$ est apprise uniquement à partir des molécules qui portent ou non le descripteur $j$.
Les descripteurs rares (31 exemples pour le plus rare) apprennent donc mal, alors que les confusions du modèle suivent les co-occurrences entre descripteurs \cite{Han-2026}.
Nous proposons de contraindre la matrice sous la forme
\begin{equation}
   \label{eq:Factorisation}
   \mathbf{W} = \mathbf{E} \, \mathbf{P},
   \qquad
   \mathbf{E} \in \mathbb{R}^{138 \times d},
   \quad
   \mathbf{P} \in \mathbb{R}^{d \times 256},
   \quad
   d < 138
\end{equation}
où $\mathbf{E}$ est figée et décrit chaque descripteur, par exemple par une décomposition en valeurs singulières de la matrice de co-occurrence calculée sur l'ensemble d'entraînement, et où seule $\mathbf{P}$ est apprise.
Chaque $\mathbf{w}_j = \mathbf{e}_j \mathbf{P}$ est alors une combinaison de $d$ vecteurs partagés~: deux descripteurs aux profils de co-occurrence proches ont nécessairement des vecteurs proches, et les descripteurs rares bénéficient des exemples de leurs voisins.
La condition $d < 138$ est essentielle~: si $d = 138$ et $\mathbf{E}$ est inversible, tout $\mathbf{W}$ reste atteignable et la contrainte disparaît.
La dimension $d$ règle donc le degré de partage imposé entre descripteurs.
Cette approche est un premier pas vers une représentation conjointe des molécules et des descripteurs, qui permettrait aussi de prédire des descripteurs absents de l'entraînement.

\paragraph{Agréger des données, notamment stéréochimiques.}
GS-LF est déjà la fusion des catalogues GoodScents et Leffingwell~; récupérer leurs versions brutes, par exemple via la plateforme Pyrfume, apporterait peu de nouvelles molécules.
Leur intérêt est ailleurs~: GS-LF a perdu la stéréochimie lors de sa constitution (les SMILES sont non stéréospécifiques), et les sources d'origine pourraient la contenir, ainsi qu'un vocabulaire plus riche.
D'autres bases (Dravnieks, Keller, Arctander) apporteraient de nouvelles molécules, avec leurs propres vocabulaires.
L'agrégation impose d'harmoniser ces vocabulaires, de résoudre les annotations contradictoires et, surtout, d'exclure toute molécule du test de GS-LF pour préserver la comparabilité.
La factorisation~\eqref{eq:Factorisation} s'y prête naturellement~: chaque source peut avoir ses propres lignes dans $\mathbf{E}$ tout en partageant $\mathbf{P}$.

L'objectif est de donner un réel avantage à l'encodeur 3D par rapport à un encodeur topologique.
Mais une limite doit être levée~: Uni-Mol ne voit que des distances, qui sont identiques pour deux énantiomères.
Les données chirales n'aideront donc que pour les diastéréoisomères, sauf à fournir explicitement la configuration des centres stéréogènes au modèle, par exemple comme caractéristique de chaque atome.
Un test simple, comparer les sorties du modèle pour les deux énantiomères de la carvone, permettra de vérifier ce point avant d'engager l'agrégation.

\paragraph{Expliquer les prédictions.}
Le sujet demande d'expliquer le fonctionnement du modèle sur quelques molécules.
Les cartes d'attention d'Uni-Mol et des méthodes d'attribution par atome permettraient de vérifier si les prédictions s'appuient sur des groupes fonctionnels connus pour leur rôle olfactif (esters pour le fruité, composés soufrés, etc.).

\paragraph{Compléter l'état de l'art.}
Les travaux GraphNOSE \cite{Sharma-2026} et sur les représentations à double graphe \cite{Zhang-2026}, identifiés en fin de phase, seront étudiés et comparés à notre approche.
